\section*{Abstract}

Der barrierefreie Zugang zu Informationen in Leichter Sprache ist eine grundlegende Voraussetzung für gesellschaftliche Teilhabe, Chancengleichheit und selbstbestimmte Information. Vor dem Hintergrund gesetzlicher Vorgaben wie dem Barrierefreiheitsstärkungsgesetz (BFSG) stehen Institutionen und Unternehmen vor der Herausforderung, komplexe Alltagstexte in großem Umfang verlässlich zu vereinfachen. Bisherige datengetriebene Ansätze zur maschinellen Textvereinfachung (Automatic Text Simplification, ATS) für das Deutsche scheiterten jedoch an drei zentralen Hürden: einem akuten Mangel an qualitätsgesicherten parallelen Trainingskorpora, dem Fehlen stufenloser und differenzierbarer Evaluierungsmetriken für Sprachkomplexität sowie der unzureichenden linguistischen Regeleinhaltung generativer Sprachmodelle.

Die vorliegende Masterarbeit begegnet diesen Herausforderungen durch die Konzeption, Implementierung und empirische Evaluation eines in sich geschlossenen, modularen Frameworks für die deutsche Leichte Sprache.

Im ersten Schritt wurde eine automatisierte Pipeline zur Korpusgewinnung entwickelt. Durch die gezielte Extraktion der Hauptinhalte (DOM-Content-Isolation), regelbasierte Bereinigungsschritte und eine semantische Ähnlichkeitsfilterung im Intervall $[0{,}60; 0{,}98]$ über Text-Embeddings konnte aus $12$ heterogenen Webquellen ein bereinigtes Parallelkorpus von $860$ konsistenten Artikelpaaren aufgebaut werden, das störende Web-Artefakte und Formatierungsreste verlässlich entfernt.

Zur referenzfreien Bewertung der sprachlichen Einfachheit wurde im zweiten Schritt ein kontinuierlicher Komplexitätsregressor auf Basis eines bidirektionalen LSTMs (BiLSTM) konzipiert. Durch das kontrollierte Mischen latenter Satzrepräsentationen (Latent MixUp) erlernt das Modell kontinuierliche Zwischenstufen der Textschwere, ohne auf künstliche Zwischentexte angewiesen zu sein. Auf einem unabhängigen Testset der \textit{Lebenshilfe Schleswig-Holstein} erzielt die Metrik eine Paar-Rangtreue von knapp $80\,\%$, korreliert stark mit qualitativen Expertenurteilen ($\rho = 0{,}6750$) und generalisiert stabil auf externe Sprachkorpora (MERLIN und TextComplexityDE).

Die generative Übersetzung wird im dritten Schritt als Sequenz-zu-Sequenz-Transformation (mBART-50) mit parametereffizienten LoRA-Adaptern modelliert. Während das überwachte Feintuning (SFT) die grundlegende Textverkürzung etabliert (Satzlänge $12{,}50 \to 7{,}61$ Wörter), schärft die anschließende direkte Präferenzoptimierung (DPO) anhand des gelernten MixUp-Belohnungssignals die formale Regeleinhaltung gezielt nach und reduziert Passiv- sowie Genitivstrukturen auf unter $1\,\%$.

In einer verblindeten Expertenstudie mit einer zertifizierten Fachkraft der Lebenshilfe bestätigt sich eine kontinuierliche Qualitätssteigerung über alle Modellstufen hinweg ($1{,}58 \to 3{,}00 \to 3{,}50 \to 3{,}92$). Eine ergänzende Konsistenzanalyse zeigt jedoch auf, dass gängige Ähnlichkeitsmetriken bei isolierten Faktenfehlern an Grenzen stoßen, womit die weitere Qualitätssteigerung des Trainingsalignments sowie NLI-basierte Konsistenzprüfungen als zentrale Hebel für zukünftige Arbeiten identifiziert werden. Die Arbeit liefert damit den methodischen Nachweis, dass eine autonome, referenzfreie Übersetzungspipeline für deutsche Leichte Sprache prinzipiell realisierbar ist.
